\section{Quantum Rejection Sampling}\label{sec:sampling}

We present the Quantum Rejection Sampling method in two steps:
\begin{itemize}
    \item {\bf Ancilla augmentation:} Find an ancilla augmentation of a distribution, such that there is a unitary $\unitarysymbol$ preparing a state with the augmented distribution as measurement distribution
    \item {\bf Acceptance amplification:} Amplify the amplitude of the ancilla qubits in state $1$ by repeated application of the Grover operator
\end{itemize}
Compared with the rejection sampling approach on the start augmentation of the distribution, we achieve a square-root speed-up in the expected number of gates to be applied per sample.

We in particular apply this scheme to sample from \ComputationActivationNetworks{}, where we apply the q-sample preparation scheme of \secref{sec:canQSample}.


\subsection{Preparation of q-samples to \CompActNets{}}\label{sec:canQSample}

%\subsection{Preparing ancilla augmented distributions for \ComputationActivationNetworks{}}

We investigate, how the above circuit encoding schemes can be applied in the preparation of q-samples to \ComputationActivationNetworks{}.
In particular, we build compositions of \computationCircuits{} and \activationCircuits{}, for which conditional distributions coincide with \ComputationActivationNetworks{}.
\red{Note that the resulting circuits are graph-controlled circuits, equivalent to Bayesian networks.}

The computation network consists already of directed cores and therefore is already a Bayesian network.
The activation network however needs ancilla augmentation.

\begin{figure}[t]
    \begin{center}
        \stantikzpic{
    \drawtextnode{-2}{4}{$a)$}

    \drawsmalltensor{1}{4}{$\acttensorof{0}$}
    \drawsmalltensor{9}{4}{$\acttensorof{\seldim\shortminus1}$}

    \drawvariabledot{1}{2.5}
    \draw[->-] (1,1) -- (1,3) node[midway,left] {\colorlabelsize$\headvariableof{0}$};
    \node[anchor=center] at (5,2) {$\cdots$};
    \drawvariabledot{9}{2.5}
    \draw[->-] (9,1) -- (9,3) node[midway,right] {\colorlabelsize$\headvariableof{\seldim\shortminus1}$};
    \drawtensorblock{5}{0}{5}{$\bencodingof{\sstat}$}

    \draw[->-] (1,-3) -- (1,-1) node[midway,left] {\colorlabelsize$\catvariableof{0}$};
    \node[anchor=center] at (5,-2) {$\cdots$};
    \draw[->-] (9,-3) -- (9,-1) node[midway,right] {\colorlabelsize$\catvariableof{\catorder\shortminus1}$};


    \begin{scope}[shift={(20,0)}]
        \drawtextnode{-2}{4}{$b)$}

        \drawtdwire{1}{7}{$\avariableof{0}$}
        \drawsmalltensor{1}{4}{$\anaugmentationof{0}$}
        \drawtdwire{9}{7}{$\avariableof{\seldim\shortminus1}$}
        \drawsmalltensor{9}{4}{$\anaugmentationof{\seldim\shortminus1}$}

        \draw[->-] (1,1) -- (1,3) node[midway,left] {\colorlabelsize$\headvariableof{0}$};
        \node[anchor=center] at (5,2) {$\cdots$};
        \draw[->-] (9,1) -- (9,3) node[midway,right] {\colorlabelsize$\headvariableof{\seldim\shortminus1}$};
        \drawtensorblock{5}{0}{5}{$\bencodingof{\sstat}$}

        \draw[->-] (1,-3) -- (1,-1) node[midway,left] {\colorlabelsize$\catvariableof{0}$};
        \node[anchor=center] at (5,-2) {$\cdots$};
        \draw[->-] (9,-3) -- (9,-1) node[midway,right] {\colorlabelsize$\catvariableof{\catorder\shortminus1}$};
    \end{scope}

}
    \end{center}
    \caption{Ancilla augmentation of an elementary \ComputationActivationNetwork{} a) by replacing each leg vector of the elementary activation tensor by an directed ancilla augmentation.}
    \label{fig:canElementaryAugmentation}
\end{figure}

%% Elementary CAN Augmentation
Elementary \ComputationActivationNetworks{} can be augmented by ancilla augmentation of each leg vector in the activation tensor network (see \figref{fig:canElementaryAugmentation}).
We understand ancilla variables as variables in a Bayesian Network having single parents, corresponding with the ancilla augmentation of an elementary tensor.
For an example, see \figref{fig:toyAccounting}.

%% Hidden variables in activation graph
To a generic activation hypergraph $\graph$, we would do ancilla augmentation for the activation tensor network, where each hidden variable is treated as a distributed qubit.
To treat it as a distributed qubit, each hidden activation qubit is prepared in uniform state (i.e. a Hadamard gate acting on a ground state) and will be omitted in the measurement scheme (either not measured, or its measurement ignored).
This can be understood from a tensor network perspective, where marginalization means contracting, which is exactly how hidden activation variables are used.
For an example in the $\ttformat$ format, see \figref{fig:canTTAugmentation}.

\begin{figure}[t]
    \begin{center}
        \begin{tikzpicture}[scale=0.35,thick] % , baseline = -3.5pt

    \node[anchor=center] (text) at (-2,4) {${a)}$};

    \drawsmalltensor{1}{5}{$\acttensorof{0}$}
    \drawsmalltensor{4}{5}{$\acttensorof{1}$}
    \drawsmalltensor{9}{5}{$\acttensorof{\seldim\shortminus1}$}

    \drawvariabledot{2.5}{3}
    \draw (2.5,1.5) -- (2.5,3) node[midway,left] {\colorlabelsize$\thirdcatvariableof{0}$};

    \draw (1.5,4) to[bend right = 50] (2.5,3);
    \draw (3.5,4) to[bend right = -50] (2.5,3);

    \draw (4.5,4) to[bend right = 50] (5.5,3);
    \drawvariabledot{5.5}{3}
    \draw (5.5,1.5) -- (5.5,3) node[midway,left] {\colorlabelsize$\thirdcatvariableof{1}$};

    \draw (8.5,4) to[bend right = -50] (7.5,3);
    \drawvariabledot{7.5}{3}
    \draw (7.5,1.5) -- (7.5,3) node[midway,right] {\colorlabelsize $\thirdcatvariableof{\seldim\shortminus2}$};

    \draw (0.5,3) -- (0.5,4);
    \drawvariabledot{0.5}{2.5}
    \draw[->-] (0.5,1) -- (0.5,3) node[midway,left] {\colorlabelsize$\headvariableof{0}$};

    \draw (4,3) -- (4,4);
    \drawvariabledot{4}{2.5}
    \draw[->-] (4,1) -- (4,3) node[midway,left] {\colorlabelsize$\headvariableof{1}$};

    \node[anchor=center] at (6.5,3) {\colorlabelsize $\cdots$};
    \draw (9.5,3) -- (9.5,4);
    \drawvariabledot{9.5}{2.5}
    \draw[->-] (9.5,1) -- (9.5,3) node[midway,right] {\colorlabelsize$\headvariableof{\seldim\shortminus1}$};
    \drawtensorblock{5}{0}{5}{$\bencodingof{\sstat}$}

    \draw[->-] (1,-3) -- (1,-1) node[midway,left] {\colorlabelsize$\catvariableof{0}$};
    \node[anchor=center] at (5,-2) {$\cdots$};
    \draw[->-] (9,-3) -- (9,-1) node[midway,right] {\colorlabelsize$\catvariableof{\catorder\shortminus1}$};


    \begin{scope}[shift={(20,0)}]
        \node[anchor=center] (text) at (-2,4) {${b)}$};

        \drawtdwire{1}{8}{$\avariableof{0}$}
        \drawsmalltensor{1}{5}{$\anaugmentationof{0}$}
        \drawtdwire{4}{8}{$\avariableof{1}$}
        \drawsmalltensor{4}{5}{$\anaugmentationof{1}$}
        \drawtdwire{9}{8}{$\avariableof{\seldim\shortminus1}$}
        \drawsmalltensor{9}{5}{$\anaugmentationof{\seldim\shortminus1}$}

        \drawvariabledot{2.5}{3}
        \draw (2.5,1.5) -- (2.5,3) node[midway,left] {\colorlabelsize$\thirdcatvariableof{0}$};

        \draw[-<-] (1.5,4) to[bend right = 50] (2.5,3);
        \draw[-<-] (3.5,4) to[bend right = -50] (2.5,3);

        \draw[-<-] (4.5,4) to[bend right = 50] (5.5,3);
        \drawvariabledot{5.5}{3}
        \draw (5.5,1.5) -- (5.5,3) node[midway,left] {\colorlabelsize$\thirdcatvariableof{1}$};

        \draw[-<-] (8.5,4) to[bend right = -50] (7.5,3);
        \drawvariabledot{7.5}{3}
        \draw (7.5,1.5) -- (7.5,3) node[midway,right] {\colorlabelsize $\thirdcatvariableof{\seldim\shortminus2}$};

        \draw[->-] (0.5,3) -- (0.5,4);
        \drawvariabledot{0.5}{2.5}
        \draw[->-] (0.5,1) -- (0.5,3) node[midway,left] {\colorlabelsize$\headvariableof{0}$};

        \draw (4,3) -- (4,4);
        \drawvariabledot{4}{2.5}
        \draw[->-] (4,1) -- (4,3) node[midway,left] {\colorlabelsize$\headvariableof{1}$};

        \node[anchor=center] at (6.5,3) {\colorlabelsize $\cdots$};
        \draw[->-] (9.5,3) -- (9.5,4);
        \drawvariabledot{9.5}{2.5}
        \draw[->-] (9.5,1) -- (9.5,3) node[midway,right] {\colorlabelsize$\headvariableof{\seldim\shortminus1}$};
        \drawtensorblock{5}{0}{5}{$\bencodingof{\sstat}$}

        \draw[->-] (1,-3) -- (1,-1) node[midway,left] {\colorlabelsize$\catvariableof{0}$};
        \node[anchor=center] at (5,-2) {$\cdots$};
        \draw[->-] (9,-3) -- (9,-1) node[midway,right] {\colorlabelsize$\catvariableof{\catorder\shortminus1}$};

    \end{scope}

\end{tikzpicture}
    \end{center}
    \caption{Ancilla augmentation of an $\ttformat$ \ComputationActivationNetwork{} a) by a Bayesian Network b).
    The hidden variables in the $\ttformat$ activation tensor are treated as distributed variables.
    Marginalization over the hidden variables and conditioning on the activation variables in state $1$, we reproduce the \ComputationActivationNetwork{}.}
    \label{fig:canTTAugmentation}
\end{figure}

\subsubsection{Sampling from \ComputationActivationNetworks{}}

For \ComputationActivationNetworks{} we find a q-sample by concatenating \computationCircuits{} and \activationCircuits{} (see \figref{fig:caCircuit}).

\begin{figure}
    \begin{center}
        \stantikzpic{

    %% Ancilla Qubits
    \drawsmalltensor{0}{17}{$\fbasis$}
    \draw (1,17) -- (13.5,17);
    \drawtensorblock{15}{17}{1.5}{$\qcaencodingof{\hypercoreof{\seldim-1}}$}
    \draw (16.5,17) -- (20,17);

    \draw (15,10) -- (15,16);
    \drawvariabledot{15}{10}

    \node[rotate=45] at (13.5,15) {$\cdots$};
    \node[rotate=45] at (13.5,8) {$\cdots$};

    \drawtextnode{0}{15.25}{$\vdots$};
    \drawtextnode{-3.5}{15}{Ancilla\\Qubits};

    \drawsmalltensor{0}{13}{$\fbasis$}
    \draw (1,13) -- (11,13);
    \drawsmalltensor{12}{13}{$\qcaencodingof{\hypercoreof{0}}$}
    \draw (13,13) -- (20,13);

    \draw (12,6) -- (12,12);
    \drawvariabledot{12}{6}

    %% Statistic Qubits
    \draw[dashed] (-3,11.5) -- (26.5,11.5);
    \drawsmalltensor{0}{10}{$\fbasis$}
    \draw (1,10) -- (7.5,10);
    \draw (7.5,9) rectangle (10.5,11);
    \drawtextnode{9}{10}{$\qcbencodingof{f_{\seldim-1}}$}
    \draw (10.5,10) -- (20,10);

    \draw (8.5,3) -- (8.5,9);
    \drawvariabledot{8.5}{3}
    \draw (9.5,-1) -- (9.5,9);
    \drawvariabledot{9.5}{-1}

    \node[rotate=45] at (7.5,8) {$\cdots$};
    \node[rotate=45] at (7.5,1) {$\cdots$};

    \drawtextnode{0}{8.25}{$\vdots$};
    \drawtextnode{-3.5}{8}{Statistic\\Qubits};

    \drawsmalltensor{0}{6}{\corelabelsize $\fbasis$}
    \draw (1,6) -- (5,6);
    \drawsmalltensor{6}{6}{$\qcbencodingof{f_{0}}$}
    \draw (7,6) -- (20,6);

    \draw (5.5,3) -- (5.5,5);
    \drawvariabledot{5.5}{3}
    \draw (6.5,-1) -- (6.5,5);
    \drawvariabledot{6.5}{-1}

    %% Distributed Qubits
    \draw[dashed] (-3,4.5) -- (26.5,4.5);

    \drawsmalltensor{0}{3}{$\fbasis$}
    \draw (1,3) -- (2,3);
    \drawsmalltensor{3}{3}{$\hgate$}

    \draw (4,3) -- (20,3);

    \drawtextnode{0}{1.25}{$\vdots$};
    \drawtextnode{-3.5}{1}{Distributed\\Qubits};

    \drawsmalltensor{0}{-1}{$\fbasis$}
    \draw (1,-1) -- (2,-1);
    \drawsmalltensor{3}{-1}{$\hgate$}

    \draw (4,-1) -- (20,-1);

    %% Measurements
    \drawqcmeasuresymbol{21}{17}
    \drawhwire{22}{17}{$\avariableof{\seldim\shortminus1}$}
    \drawsmalltensor{25}{17}{$\tbasis$}

    \drawtextnode{23}{15.5}{$\vdots$};
    \drawqcmeasuresymbol{21}{13}
    \drawhwire{22}{13}{$\avariableof{0}$};
    \drawsmalltensor{25}{13}{$\tbasis$}

    \drawqcmeasuresymbol{21}{10}
    \drawhwire{22}{10}{$\headvariableof{\seldim\shortminus1}$}
    \drawvariabledot{24}{10}
    \drawtextnode{23}{8.5}{$\vdots$};
    \drawqcmeasuresymbol{21}{6}
    \drawhwire{22}{6}{$\headvariableof{0}$}
    \drawvariabledot{24}{6}

    \drawqcmeasuresymbol{21}{3}
    \drawhwire{22}{3}{$\catvariableof{\catorder\shortminus1}$}
    \drawtextnode{23}{1.5}{$\vdots$};
    \drawqcmeasuresymbol{21}{-1}
    \drawhwire{22}{-1}{$\catvariableof{0}$}

    \begin{scope}[shift={(-15,0)}]
        \drawhwire{3}{3}{$\catvariableof{\catorder\shortminus1}$}
        \drawhwire{3}{-1}{$\catvariableof{0}$}
        \drawtextnode{4}{1.5}{$\vdots$}
        \draw (-1,4) rectangle (3,-2);
        \drawtextnode{1}{1}{$\probof{\hybridparam}$};
        \drawtextnode{-2}{1}{$\acceptanceprob\cdot$}
        \drawtextnode{6}{1}{${=}$}
    \end{scope}

}
    \end{center}
    \caption{
        Quantum Circuit to reproduce a q-sample of an augmented \ComputationActivationNetwork{}.
        We measure the distributed qubits $\shortcatvariables$ and the ancilla qubits $\avariableof{[\seldim]}$ and reject all samples, where an ancilla qubit is measured as $0$.
    }\label{fig:caCircuit}
\end{figure}

\begin{example}[Toy Accounting Example]
    \label{exa:toyAccounting}
    For a more detailed example, let us consider a system of three variables $A1$ Account 1 is booked, $A2$ Account 2 is booked, $F$ a feature on an invoice.
    Assume the following two rules have to be respected:
    \begin{itemize}
        \item \textcolor{\concolor}{Exactly one account must be booked.}
        \item \textcolor{\probcolor}{If feature $\mathrm{F}$ is present on the invoice, the account $\mathrm{A1}$ is typically booked.}
    \end{itemize}
    We formalize this with the statistic
    \begin{align*}
        \hlnstat = (\catvariableof{A1} \oplus \catvariableof{A2}, \catvariableof{F}\Rightarrow \catvariableof{A1})\, .
    \end{align*}
    Any elementary \ComputationActivationNetwork{} with the statistic $\hlnstat$ can be realized by the circuit shown in \figref{fig:toyAccounting}.
    %In particular, the \computationCircuit{} consists of:
    For the preparation of the first statistic qubit $\headvariableof{0}$ to $\sstatcoordinateof{0}=\catvariableof{A1} \oplus \catvariableof{A2}$ by a \computationCircuit{} we exploit that
    \begin{align*}
        \catvariableof{A1}\oplus\catvariableof{A2}
        &= \onehotmapofat{1}{\catvariableof{A1}} \otimes \onehotmapofat{0}{\catvariableof{A2}}
        + \onehotmapofat{0}{\catvariableof{A1}} \otimes \onehotmapofat{1}{\catvariableof{A2}} \\
        &= \left(\onehotmapofat{1}{\catvariableof{A1}} \otimes \onehotmapofat{0}{\catvariableof{A2}}\right)
        \oplus \left(\onehotmapofat{0}{\catvariableof{A1}} \otimes \onehotmapofat{1}{\catvariableof{A2}}\right)
    \end{align*}
    where by $\oplus$ we denote coordinatewise summation mod 2.
    The statistic qubit $\headvariableof{0}$ is therefore prepared by two controlled NOT gates (see \figref{fig:toyAccounting}b).

    For the preparation of the second statistic qubit $\headvariableof{1}$ to $\sstatcoordinateof{1}=\catvariableof{F}\Rightarrow \catvariableof{A1}$ we exploit that the implication is $\truesymbol$ except for the case where the premise $\catvariableof{F}$ is $\truesymbol$ and the head $\catvariableof{A1}$ is $\falsesymbol$.
    In out mod 2 calculus, this amounts to
    \begin{align*}
        \catvariableof{F}\Rightarrow \catvariableof{A1}
        = \onesat{\catvariableof{F},\catvariableof{A1}} \oplus \onehotmapofat{1}{\catvariableof{F}} \otimes \onehotmapofat{0}{\catvariableof{A1}}\, .
    \end{align*}
    It follows that the statistic qubit $\headvariableof{1}$ is prepared by a NOT gate (that is a Pauli-X gate $\paulixsymbol$) and a second controlled NOT gate (see \figref{fig:toyAccounting}b).

    For both cases, the preparation of the statistic qubit can be done by two controlled NOT gates exploiting the \polynomialSparsity{} of the connectives.

    Any non-elementary \ComputationActivationNetwork{} with the statistic $\hlnstat$ can be prepared by the circuit, when choosing a single ancilla qubit, which is uniformly controlled by the qubits $\headvariableof{[2]}$.

    \begin{figure}[t]
        \stantikzpic{
            \drawtextnode{-16}{14}{$a)$}
            \drawtextnode{-4}{14}{$b)$}

            \begin{scope}[shift={(-12,2.25)}]

                \drawvariablenode{-1.5}{10}{$\avariableof{0}$}{A0}
                \drawvariablenode{1.5}{10}{$\avariableof{1}$}{A1}

                \drawvariablenode{-1.5}{5}{$\headvariableof{0}$}{H0}
                \drawvariablenode{1.5}{5}{$\headvariableof{1}$}{H1}

                \drawvariablenode{-3}{0}{$\catvariableof{A2}$}{XA2}
                \drawvariablenode{0}{0}{$\catvariableof{A1}$}{XA1}
                \drawvariablenode{3}{0}{$\catvariableof{F}$}{F}

                \draw[->-] (H0) -- (A0);
                \draw[->-] (H1) -- (A1);

                \coordinate (F0) at (-1.5,2.5);
                \draw[->-] (XA2) -- (F0);
                \draw[->-] (XA1) -- (F0);
                \draw[->-] (F0) -- (H0);

                \coordinate (F1) at (1.5,2.5);
                \draw[->-] (F) -- (F1);
                \draw[->-] (XA1) -- (F1);
                \draw[->-] (F1) -- (H1);
            \end{scope}

        %% Ancilla Qubits

            \drawsmalltensor{0}{14}{$\onehotmapof{0}$}
            \draw (1,14) -- (13.5,14);
            \drawsmalltensor{14.5}{14}{$\qcaencodingof{\hypercoreof{1}}$}
            \draw (15.5,14) -- (18,14);
            \drawqcmeasuresymbol{19}{14}
            \draw (20,14) -- (21.5,14) node[above,midway] {\colorlabelsize $\avariableof{1}$};

            \draw (14.5,8.5) -- (14.5,13);
            \drawvariabledot{14.5}{8.5}

            \drawsmalltensor{0}{11.5}{$\onehotmapof{0}$}
            \draw (1,11.5) -- (11.5,11.5);
            \drawsmalltensor{12.5}{11.5}{$\qcaencodingof{\hypercoreof{0}}$}
            \draw (13.5,11.5) -- (18,11.5);
            \drawqcmeasuresymbol{19}{11.5}
            \draw (20,11.5) -- (21.5,11.5) node[above,midway] {\colorlabelsize $\avariableof{0}$};

            \draw (12.5,6) -- (12.5,10.5);
            \drawvariabledot{12.5}{6}

            %% Statistic Qubits
            \draw[dashed] (-3,10) -- (21,10);

            \drawsmalltensor{0}{8.5}{$\fbasis$}
            \draw (1,8.5) -- (4,8.5);
            \drawsmalltensor{5}{8.5}{$\paulixsymbol$}
            \draw (6,8.5) -- (18,8.5);
            \drawcnotsymbol{11}{8.5}
            \drawqcmeasuresymbol{19}{8.5}
            \draw (20,8.5) -- (21.5,8.5) node[above,midway] {\colorlabelsize $\headvariableof{1}$};
            \drawvariabledot{21.5}{8.5}

            \drawsmalltensor{0}{6}{$\fbasis$}
            \draw (1,6) -- (18,6);
            \drawcnotsymbol{9}{6}
            \drawqcmeasuresymbol{19}{6}
            \draw (20,6) -- (21.5,6) node[above,midway] {\colorlabelsize $\headvariableof{0}$};
            \drawvariabledot{21.5}{6}

            \drawcnotsymbol{5}{6}

            %% Distributed Qubits
            \draw[dashed] (-3,4.5) -- (21,4.5);

            \drawsmalltensor{0}{3}{$\fbasis$}
            \draw (1,3) -- (2,3);
            \drawsmalltensor{3}{3}{$\hgate$}
            \draw (4,3) -- (18,3);
            \drawqcmeasuresymbol{19}{3}
            \draw (20,3) -- (21.5,3) node[above,midway] {\colorlabelsize $\catvariableof{F}$};

            \drawsmalltensor{0}{0.5}{$\fbasis$}
            \draw (1,0.5) -- (2,0.5);
            \drawsmalltensor{3}{0.5}{$\hgate$}
            \draw (4,0.5) -- (6,0.5);
            \drawsmalltensor{7}{0.5}{$\paulixsymbol$}

            \draw (8,0.5) -- (18,0.5);
            \drawqcmeasuresymbol{19}{0.5}
            \draw (20,0.5) -- (21.5,0.5) node[above,midway] {\colorlabelsize $\catvariableof{A2}$};

            \drawsmalltensor{0}{-2}{$\fbasis$}
            \draw (1,-2) -- (2,-2);
            \drawsmalltensor{3}{-2}{$\hgate$}
            \draw (4,-2) -- (6,-2);
            \drawsmalltensor{7}{-2}{$\paulixsymbol$}
            \draw (8,-2) -- (12,-2);
            \drawsmalltensor{13}{-2}{$\paulixsymbol$}
            \draw (14,-2) -- (18,-2);
            \drawqcmeasuresymbol{19}{-2}
            \draw (20,-2) -- (21.5,-2) node[above,midway] {\colorlabelsize $\catvariableof{A1}$};

            \draw (5.25,5) -- (5.25,-2);
            \draw (4.75,5) -- (4.75,0.5);

            \drawvariabledot{4.75}{0.5}
            \drawvariabledot{5.25}{-2}

            \draw (9.25,5) -- (9.25,-2);
            \draw (8.75,5) -- (8.75,0.5);

            \drawvariabledot{8.75}{0.5}
            \drawvariabledot{9.25}{-2}

            \draw (11.25,7.5) -- (11.25,-2);
            \draw (10.75,7.5) -- (10.75,3);
            \drawvariabledot{10.75}{3}
            \drawvariabledot{11.25}{-2}
        }
        \caption{Circuit to sample from the \ComputationActivationNetwork{} in the toy accounting \exaref{exa:toyAccounting}.
        %Note that a first Pauli-X gate $\paulixsymbol$ acting on $\catvariableof{A2}$ can be omitted, since the prepared stated is uniform.
        }
        \label{fig:toyAccounting}
    \end{figure}

\end{example}


\subsubsection{Acceptance Amplification}

Let there be a quantum state, which measurement distribution is an ancilla augmentation $\secprobat{\avariables,\shortcatvariables}$ of a target distribution $\probwith$.
We now show how Grover operators can be used to amplify the acceptance state and improve the efficiency of the rejection sampling method through an increase of the acceptance probability.
This is an instance of amplitude amplification.

\subsection{Grover construction}

We relate the state $\qstateof{\goodstatesymbol}$ with the event $\avariables=\onetuple{[\seldim]}$, by constructing it as
\begin{align*}
    \qstateofat{\goodstatesymbol}{\avariables,\headvariables,\shortcatvariables}
    =\onehotmapofat{\onetuple{[\seldim]}}{\avariables}
    \otimes \contractionof{\bencodingofat{\sstat}{\headvariables,\shortcatvariables},\sqrt{\probwith}}{\headvariables,\shortcatvariables}
\end{align*}
The complement state $\qstateof{\badstatesymbol}$ is constructed if $\acceptanceprob\neq1$ (in which case the acceptance amplilification scheme would not be necessary)
\begin{align*}
    \probofat{\badstatesymbol}{\avariables,\headvariables,\shortcatvariables}
    = \frac{1}{1-\acceptanceprob}
    \cdot\left(
             \secprobat{\avariables,\headvariables,\shortcatvariables}-\acceptanceprob \cdot \onehotmapofat{\onetuple{[\seldim]}}{\avariables}\otimes\probat{\headvariables,\shortcatvariables}
    \right)
\end{align*}
and
\begin{align*}
    \qstateofat{\badstatesymbol}{\avariables,\headvariables,\shortcatvariables}=\sqrt{\probofat{\perp}{\avariables,\headvariables,\shortcatvariables}} \, .
\end{align*}

% Grover operator instance
%We build the Grover operator (see \secref{sec:groverOperator}) with respect to
%\begin{itemize}
%%    \item $$
%    \item $\qstateof{\badstatesymbol}=\sqrt{\probof{\perp}}$
%\end{itemize}
%and the state

We thus have
\begin{align*}
    \qstateofat{0}{\avariables,\headvariables,\shortcatvariables}
    = \sinof{\frac{\rotanglesymbol}{2}} \qstateofat{0}{\avariables,\headvariables,\shortcatvariables}
    + \cosof{\frac{\rotanglesymbol}{2}} \qstateofat{1}{\avariables,\headvariables,\shortcatvariables}
\end{align*}
where the $\rotanglesymbol$ is related to the acceptance probability as
\begin{align*}
    \sinof{\frac{\rotanglesymbol}{2}}
    = \sqrt{\acceptanceprob} \, .
\end{align*}



\subsubsection{Circuit Realization}

Both reflection operators are realized as their negations using the Pauli-Z scheme (see \secref{sec:groverOperator}).
The negations chancel in their concatenation, which is therefore the Grover operator.


\subsubsection{Ancilla augmentation preserved}

The measurement distribution to any quantum state in $\subspaceof{}\backslash\spanof{\qstateof{\badstatesymbol}}$ is an ancilla augmentation of $\probwith$. % Need to show?
When amplifying the acceptance probability we stay in the subspace $\subspaceof{}$.
When we do not overrotate and since we start with $\acceptanceprob>0$ the amplified state is never in $\spanof{\qstateof{\badstatesymbol}}$.

%% Realization of the grover operator
%While the reflection $\reflectionof{0}$ is realized as described in \secref{sec:groverOperator}, the reflection $\reflectionof{\badstatesymbol}$ is a multi-controlled Pauli-Z on the ancilla variables. \
%\red{No, we use the negated about the good state!}



\subsection{Optimal rotation numbers}

% Application
Now the application of $\repnum$ grover operators $\groversymbol$ on $\qstateof{0}$ prepares a state
\begin{align*}
    \qstateofat{\repnum}{\avariables,\shortcatvariables}
    = \sinof{\left(\frac{1}{2}+\repnum\right)\rotanglesymbol} \qstateofat{\goodstatesymbol}{\avariables,\shortcatvariables}
    + \cosof{\left(\frac{1}{2}+\repnum\right)\rotanglesymbol} \qstateofat{\badstatesymbol}{\avariables,\shortcatvariables}
\end{align*}
The state is closest to $\qstateof{\goodstatesymbol}$ when
\begin{align*}
    \left(\frac{1}{2}+\repnum\right)
    \rotanglesymbol \approx \frac{\pi}{2} \, .
\end{align*}

Estimating for small $\acceptanceprob$ (tight for small $\acceptanceprob$)
\begin{align*}
    \sqrt{\acceptanceprob} = \sinof{\frac{\rotanglesymbol}{2}}\leq \frac{\rotanglesymbol}{2}
\end{align*}
we get a bound
\begin{align*}
    \repnumof{\mathrm{opt}} \leq \frac{\pi}{4\sqrt{\acceptanceprob}} \, .
\end{align*}
This shows that we have a square root improvement on classical rejection sampling, which would require $\frac{1}{\acceptanceprob}$ repetitions to produce a sample.
Note however that we compare quantum circuit length with the number of classical repetitions in sampling.


%Note, that the variable qubits are uniformly distributed when only the computation circuit is applied.
%When sampling the probability distribution, we need the ancilla qubits to be in state $1$ in order for the sample to be valid.
%Any other states will have to be rejected.

%\textbf{Open Question:} Is there a way to avoid amplitude amplification and use a more direct circuit implementation of the activation network?
%- Cannot be the case, when the encoding is determined by the activation tensor alone: Needs to use the computated statistic as well.
%Negative result in \cite[Section~6.6]{nielsen_quantum_2010}: When using multiple applications of \computationCircuit{} in combination with further unitaries, need at least as many applications as with amplitude amplification.

\subsection{Applications on \CompActNets{}}

The above scheme prepares the ancilla augmented \ComputationActivationNetwork{}
\begin{align*}
    \secprobat{\avariables,\headvariables,\shortcatvariables}
    = \acceptanceprob \cdot \onehotmapofat{\onetuple{[\seldim]}}{\avariables}
    \otimes \contractionof{\bencodingofat{\sstat}{\headvariables,\shortcatvariables},\probwith}{\headvariables,\shortcatvariables}
    + (1-\acceptanceprob) \probofat{\perp}{\shortcatvariables,\headvariables,\avariables}
\end{align*}
and the prepared state is the combination of q-samples
\begin{align*}
    \qstateofat{0}{\avariables,\headvariables,\shortcatvariables}
    = \sqrt{\acceptanceprob} \cdot \onehotmapofat{\onetuple{[\seldim]}}{\avariables}
    \otimes \contractionof{\bencodingofat{\sstat}{\headvariables,\shortcatvariables},\sqrt{\probwith}}{\headvariables,\shortcatvariables}
    + \sqrt{(1-\acceptanceprob)} \sqrt{\probofat{\perp}{\shortcatvariables,\headvariables,\avariables}} \, .
\end{align*}
%Here we understand $\headvariables$ as extending $\shortcatvariables$.

Sampling from the augmented \ComputationActivationNetwork{} can be also done classically by forward sampling:
Just draw the variables from uniform, calculate the value qubit by a logical circuit inference and accept with probability by the computed value.
To achieve a quantum advantage we manipulate the circuit first by the amplitude amplification scheme and increase the acceptance probability.

\subsection{Examples}

%% To be included:

\begin{example}[Sample a model of a propositional formula]

    Consider the task of sampling from the models of a propositional formula $\formulawith$, which can be expressed as a \ComputationActivationNetwork{} as
    \begin{align*}
        \frac{1}{\contraction{\formula}} \cdot \formulawith
        = \frac{1}{\contraction{\formula}} \cdot \contractionof{\bencodingofat{\formula}{\headvariable,\shortcatvariables},\tbasisat{\headvariable}}{\shortcatvariables} \, .
    \end{align*}

    Note that in this case, $\frac{1}{2^{\catorder}} \cdot \bencodingofat{\formula}{\headvariable,\shortcatvariables}$ is an ancilla augmentation of $\frac{1}{\contraction{\formula}} \cdot \formulawith$ with $\headvariable$ as the ancilla variable, which q-sample is prepared by the oracle $\qcbencodingof{\formula}$ acting on the ground state $\onehotmapofat{0_{[\catorder+1]}}{\headvariable,\shortcatvariables}$.
    The acceptance rate is then
    \begin{align*}
        \acceptanceprob
        = \frac{1}{2^{\catorder}} \cdot \contraction{\bencodingofat{\formula}{\headvariable=1,\shortcatvariables}}
        = \frac{\contraction{\formula}}{2^{\catorder}}
    \end{align*}
    and the optimal number of rotations is
    \begin{align*}
        \repnumof{\optsymbol}
        = 2 \cdot \mathrm{arcsin}\left(\sqrt{\frac{\contraction{\formula}}{2^{\catorder}}}\right) \, .
    \end{align*}
    The example of a single rotation is sketched in \figref{fig:formulaAmplificationCircuit}.

    \begin{figure}

        \stantikzpic{

            \drawsmalltensor{0}{9}{$\fbasis$}
            \draw (1,9) -- (4,9);
            \drawhblock{5.5}{9}{1.5}{$\qcbencodingof{\exformula}$}
            \draw (5,6) -- (5,8);
            \drawvariabledot{5}{6}
            \draw (6,2) -- (6,8);
            \drawvariabledot{6}{2}


            \draw[dashed] (7.5,12) -- (7.5,-1) node[below] {\corelabelsize $\qstateofat{0}{\headvariable,\shortcatvariables}$};
            \drawtextnode{9}{11}{\corelabelsize $\reflectionof{\goodstatesymbol}$}
            \draw[dashed] (10.5,11) -- (10.5,0.5);
            \drawtextnode{19.5}{11}{\corelabelsize $\reflectionof{0}$}

            \draw[dashed] (28.5,12) -- (28.5,-1) node[below] {\corelabelsize $\qstateofat{1}{\headvariable,\shortcatvariables}$};

            \draw (7,9) -- (8,9);
            \draw (8,8) rectangle (10,10);
            \drawsmalltensor{9}{9}{\corelabelsize $\paulizsymbol$}

            \begin{scope}[shift={(7,0)}]
                \draw (3,9) -- (4,9);
                \drawhblock{5.5}{9}{1.5}{$\qcbencodingof{\exformula}$}
                \draw (5,6) -- (5,8);
                \drawvariabledot{5}{6}
                \draw (6,2) -- (6,8);
                \drawvariabledot{6}{2}

                \draw (7,9) -- (8.5,9);
                \drawsmalltensor{9.5}{9}{$\paulixsymbol$}
                \draw (10.5,9) -- (11.5,9);
                \drawsmalltensor{12.5}{9}{$\paulizsymbol$}
                \draw (13.5,9) -- (14.5,9);
                \drawsmalltensor{15.5}{9}{$\paulixsymbol$}

                \draw (12,6) -- (12,8);
                \drawvariabledot{12}{6}
                \draw (13,2) -- (13,8);
                \drawvariabledot{13}{2}

                \drawsmalltensor{7.5}{6}{$\hgate$}
                \draw (8.5,6) -- (9.5,6);
                \drawsmalltensor{10.5}{6}{$\paulixsymbol$}
                \draw (11.5,6) -- (13.5,6);
                \drawsmalltensor{14.5}{6}{$\paulixsymbol$}
                \draw (15.5,6) -- (16.5,6);
                \drawsmalltensor{17.5}{6}{$\hgate$}
                \draw (18.5,6) -- (23,6);

                \drawsmalltensor{7.5}{2}{$\hgate$}
                \draw (8.5,2) -- (9.5,2);
                \drawsmalltensor{10.5}{2}{$\paulixsymbol$}
                \draw (11.5,2) -- (13.5,2);
                \drawsmalltensor{14.5}{2}{$\paulixsymbol$}
                \draw (15.5,2) -- (16.5,2);
                \drawsmalltensor{17.5}{2}{$\hgate$}
                \draw (18.5,2) -- (23,2);

            \end{scope}

            \begin{scope}[shift={(21,0)}]
                \draw (2.5,9) -- (4,9);
                \drawhblock{5.5}{9}{1.5}{$\qcbencodingof{\exformula}$}
                \draw (7,9) -- (9,9);

                \draw (5,6) -- (5,8);
                \drawvariabledot{5}{6}
                \draw (6,2) -- (6,8);
                \drawvariabledot{6}{2}

            \end{scope}

            \drawsmalltensor{0}{6}{$\fbasis$}
            \draw (1,6) -- (2,6);
            \drawsmalltensor{3}{6}{$\hgate$}
            \draw (4,6) -- (13.5,6);

            \drawsmalltensor{0}{2}{$\fbasis$}
            \draw (1,2) -- (2,2);
            \drawsmalltensor{3}{2}{$\hgate$}
            \draw (4,2) -- (13.5,2);

            \drawtextnode{-3}{9}{\corelabelsize $\headvariable$}
            \draw[dashed] (-2,7.5) -- (30.5,7.5);
            \drawtextnode{-3}{4}{\corelabelsize $\shortcatvariables$}
            \drawtextnode{0}{4.5}{\corelabelsize $\vdots$}

        }
        \caption{Amplitude amplifying circuit for the ancilla augmentation of a function $\exformula$, for $\repnum=1$ repetitions.
        The initial state $\qstateofat{0}{\headvariable,\shortcatvariables}$ is the q-sample of the by $\headvariable$ augmented uniform distributions among the models of $\exformula$, and prepared by action of the oracle $\qcbencodingof{\exformula}$ on the ground state.
        The amplification is then done by the Grover operator, which consists in the reflection $\reflectionof{\goodstatesymbol}$ across the state $\qstateofat{\goodstatesymbol}{\headvariable,\shortcatvariables}$ followed by a reflection reflection $\reflectionof{0}$ across the state $\qstateofat{0}{\headvariable,\shortcatvariables}$.}
        \label{fig:formulaAmplificationCircuit}
    \end{figure}
\end{example}
\begin{example}[Sampling from the toy accounting system]
    For the toy accounting system we have
    \begin{align*}
        \max_{\catindexof{F},\catindexof{A1},\catindexof{A2}\in[2]}
        \probat{\indexedcatvariableof{F},\indexedcatvariableof{A1},\indexedcatvariableof{A2}}
        = \begin{cases}
              \frac{1}{3+\expof{-\canparam}} & \ifspace \canparam \geq 0 \\
              \frac{1}{3\cdot\expof{\canparam}+1} & \ifspace \canparam < 0
        \end{cases}
    \end{align*}
    and therefore the acceptance probability when sampling $(\catvariableof{F},\catvariableof{A1},\catvariableof{A2})$ from uniform by
    \begin{align*}
        \acceptanceprob
        = \frac{1}{8} \cdot \left(\max_{\catindexof{F},\catindexof{A1},\catindexof{A2}\in[2]}
        \probat{\indexedcatvariableof{F},\indexedcatvariableof{A1},\indexedcatvariableof{A2}}\right)^{-1}
        = \begin{cases}
              \frac{3+\expof{-\canparam}}{8} & \ifspace \canparam \geq 0 \\
              \frac{3\cdot\expof{\canparam}+1}{8} & \ifspace \canparam < 0
        \end{cases} \, .
    \end{align*}
    The optimal number of amplifications is then
    \begin{align*}
        \repnumof{\optsymbol}
        = \begin{cases}
              0 & \ifspace \canparam \geq \canparamof{\mathrm{critical}} \\
              1 & \ifspace \canparam < \canparamof{\mathrm{critical}}
        \end{cases}
    \end{align*}
    for
    \begin{align*}
        \canparamof{\mathrm{critical}}
        = \lnof{1-\frac{4\cosof{1}}{3}} \approx -1.27 \, .
    \end{align*}
    Note that in the limit of $\canparam\rightarrow -\infty$, we effectively have a single model and therefore a minimum of the acceptance probability among the distributions of three Boolean variables.
\end{example}


